\section{Introduction}
Consider the general equation of motion of a dynamic system under external load $\bp\left(t\right)$,
\begin{gather}
\br\left(\bu,\bv,\ba\right)=\bp\left(t\right)
\end{gather}
A practical simplification may assume the contributions of $\bu$, $\bv=\dot{\bu}$ and $\ba=\ddot{\bu}$ are separable, indicating no coupling among them, then the resistance $\br$ can be split into three components
\begin{gather}
\br=\br_{\bu}+\br_{\bv}+\br_{\ba}.
\end{gather}
Additional energy dissipating terms can be introduced,
\begin{gather}
\br_{\bu}+\br_{\bv}+\br_{\ba}+\bbf_{\bu}+\bbf_{\bv}+\bbf_{\ba}=\bp,
\end{gather}
in which $\bbf_{\baleph}$ (for $\baleph$ being one of $\bu$, $\bv$ and $\ba$) is governed by the following differential equation:
\begin{gather}
\dot{\bbf_{\baleph}}=-s_{\baleph}\bbf_{\baleph}+m_{\baleph}\by_{\baleph}
\end{gather}
where $s_{\baleph}$ and $m_{\baleph}$ are two scalar parameters that control the response, $\by_{\baleph}=\by_{\baleph}\left(\baleph\right)$ is a function of $\baleph$, various forms can be chosen. For brevity, we consider $\by_{\baleph}=\br_{\baleph}$ in this work unless otherwise stated.
The subscript $\left(\cdot\right)_{\baleph}$ may be dropped for brevity in the following discussion, providing no ambiguity would occur by doing so.
\subsection{Free Vibration}
\subsubsection{Displacement Dependent}
\paragraph{One Single Term}
Since only the displacement dependent energy dissipating force $\bbf_{\bu}$ is considered, the subscript $\left(\cdot\right)_{\bu}$ is omitted for brevity.
By selecting $\by=\mb{K}\bu$, the linear elastic system simplifies to the following.
\begin{gather}
\left\{\begin{array}{l}
\mb{M}\ddot{\bu}+\mb{K}\bu+\bbf=\mb{0},\\
\dot{\bbf}=-s\bbf+m\mb{K}\bu.
\end{array}\right.
\end{gather}
One could obtain
\begin{gather}
\mb{M}\dddot{\bu}+\mb{K}\dot{\bu}+s\mb{M}\ddot{\bu}+\left(s+m\right)\mb{K}\bu=\mb{0}.
\end{gather}
With the assist of assuming a harmonic solution $\bu\left(t\right)=\mb{U}\exp\left(i\omega{}t\right)$, the characteristic equation can be obtained as follows.
\begin{gather}
-i\omega^3\mb{M}-s\omega^2\mb{M}+i\omega\mb{K}+\left(s+m\right)\mb{K}=\mb{0}.
\end{gather}
Rearranging gives
\begin{gather}
\left(s+m+i\omega\right)\mb{K}-\left(s+i\omega\right)\omega^2\mb{M}=\mb{0}\qquad\longrightarrow\qquad
\dfrac{s+m+i\omega}{s+i\omega}\mb{K}-\omega^2\mb{M}=\mb{0}.
\end{gather}
The dynamic stiffness is then
\begin{gather}
\mb{K}^*=\dfrac{s+m+i\omega}{s+i\omega}\mb{K}=\left(1+\dfrac{m}{s+i\omega}\right)\mb{K},
\end{gather}
the corresponding storage stiffness $\mb{K}'$ and loss stiffness $\mb{K}''$ can be expressed as
\begin{gather}
\mb{K}'=\left(1+\dfrac{sm}{s^2+\omega^2}\right)\mb{K},\qquad
\mb{K}''=\dfrac{m\omega}{s^2+\omega^2}\mb{K},
\end{gather}
such that $\mb{K}^*=\mb{K}'-\mb{K}''i$.
The ratio $\eta_u$ between them is
\begin{gather}
\eta_u=\dfrac{m\omega}{s^2+sm+\omega^2}.
\end{gather}
\paragraph{Multiple Terms}
Now consider multiple energy dissipating terms $\bbf_j$, all contributing to the equation of motion, one have
\begin{gather}
\left\{\begin{array}{l}
\mb{M}\ddot{\bu}+\mb{K}\bu+\sum\bbf_j=\mb{0},\\
\dot{\bbf}_j=-s_j\bbf_j+m_j\mb{K}\bu.
\end{array}\right.
\end{gather}
Assuming harmonic solutions for both $\bu$ and $\bbf$,
\begin{gather}
\bu\left(t\right)=\mb{U}\exp\left(i\omega{}t\right),\qquad
\bbf_j\left(t\right)=\mb{F}_j\exp\left(i\omega{}t\right),
\end{gather}
the second expression can be expressed in the frequency domain as
\begin{gather}
i\omega\mb{F}_j=-s_j\mb{F}_j+m_j\mb{K}\mb{U},
\end{gather}
thus
\begin{gather}
\mb{F}_j=\dfrac{m_j}{s_j+i\omega}\mb{K}\mb{U}.
\end{gather}
The characteristic equation can be then derived as
\begin{gather}
\left(1+\sum\dfrac{m_j}{s_j+i\omega}\right)\mb{K}-\omega^2\mb{M}=\mb{0}.
\end{gather}
The dynamic stiffness is
\begin{gather}
\mb{K}^*=\left(1+\sum\dfrac{m_j}{s_j+i\omega}\right)\mb{K}
\end{gather}
with
\begin{gather}
\mb{K}'=\left(1+\sum\dfrac{s_jm_j}{s_j^2+\omega^2}\right)\mb{K},\qquad
\mb{K}''=\sum\dfrac{m_j\omega}{s_j^2+\omega^2}\mb{K}.
\end{gather}
The ratio $\eta_u$ between them is now
\begin{gather}
\eta_u=\dfrac{\sum\dfrac{m_j\omega}{s_j^2+\omega^2}}{1+\sum\dfrac{s_jm_j}{s_j^2+\omega^2}}.
\end{gather}

The multi-term formulation is a straightforward additive extension of the single-term case.
The resulting dynamic stiffness has the same form as that reported by \citet{Huang2019}, where dissipation is introduced through resistance filtering.
The foregoing derivation therefore provides an alternative differential equation based formulation.
\subsubsection{Acceleration Dependent}
We consider multiple acceleration dependent energy dissipation forces $\bbf_j$ directly.
The system is
\begin{gather}
\left\{\begin{array}{l}
\mb{M}\ddot{\bu}+\mb{K}\bu+\sum\bbf_j=\mb{0},\\
\dot{\bbf}_j=-s_j\bbf_j+m_j\mb{M}\ba.
\end{array}\right.
\end{gather}
Adopting a similar approach, the characteristic equation can be derived as
\begin{gather}
\dfrac{1}{1+\sum\dfrac{m_j}{s_j+i\omega}}\mb{K}-\omega^2\mb{M}=\mb{0}.
\end{gather}
The loss ratio $\eta_a$ satisfies the condition $\abs{\eta_a}=\abs{\eta_u}$ since the scalar coefficient in the expression of dynamic stiffness is the reciprocal of that in the previous displacement dependent case.
\section{Implementation}
To develop a general algorithm that can be universally applied to various types of systems, we turn back to the general expression with multiple displacement dependent term.
The governing system is
\begin{gather}\label{eq:eqv_sys}
\left\{
\begin{array}{l}
\dot{\bbf}_j=-s_j\bbf_j+m_j\br_{\bu},\\
\br_{\bu}+\br_{\bv}+\br_{\ba}+\sum\bbf_j=\mb{p}.
\end{array}\right.
\end{gather}

The first expression is a first-order ODE (of $\bbf_j$) that can be integrated by the (implicit) trapezoidal rule, which is second-order accurate and A-stable,
\begin{gather}
\bbf_{j,n+1}=\bbf_{j,n}+\dfrac{\Delta{}t}{2}\left(\dot{\bbf}_{j,n}+\dot{\bbf}_{j,n+1}\right).
\end{gather}
Expanding gives
\begin{gather}
\bbf_{j,n+1}=\bbf_{j,n}+\dfrac{\Delta{}t}{2}\left(\left(-s_j\bbf_{j,n}+m_j\br_{\bu,n}\right)+\left(-s_j\bbf_{j,n+1}+m_j\br_{\bu,n+1}\right)\right),
\end{gather}
one may further rearrange to obtain
\begin{gather}\label{eq:discretised_c}
\bbf_{j,n+1}=\dfrac{2-s_j\Delta{}t}{2+s_j\Delta{}t}\bbf_{j,n}+\dfrac{m_j\Delta{}t}{2+s_j\Delta{}t}\br_{\bu,n}+\dfrac{m_j\Delta{}t}{2+s_j\Delta{}t}\br_{\bu,n+1},
\end{gather}
in which subscripts $\left(\cdot\right)_n$ and $\left(\cdot\right)_{n+1}$ denote the corresponding quantity at $t_n$ and $t_{n+1}=t_n+\Delta{}t$.

Now assuming the equation of motion is satisfied at $t_{n+1}$, the second expression is
\begin{gather}\label{eq:residual}
\br_{\bu,n+1}+\br_{\bv,n+1}+\br_{\ba,n+1}+\sum\dfrac{2-s_j\Delta{}t}{2+s_j\Delta{}t}\bbf_{j,n}+\sum\dfrac{m_j\Delta{}t}{2+s_j\Delta{}t}\br_{\bu,n}+\sum\dfrac{m_j\Delta{}t}{2+s_j\Delta{}t}\br_{\bu,n+1}=\mb{p}_{n+1}.
\end{gather}
Differentiation leads to the following revised effective stiffness $\bhat{K}_{n+1}$,
\begin{gather}\label{eq:revised_k}
\bhat{K}_{n+1}=\bbar{K}_{n+1}+\sum\dfrac{m_j\Delta{}t}{2+s_j\Delta{}t}\mb{K}_{n+1}.
\end{gather}
in which the conventional effective stiffness $\bbar{K}_{n+1}$ is
\begin{gather}
\bbar{K}_{n+1}=\mb{K}_{n+1}+\mb{C}_{n+1}\ddfrac{\bv_{n+1}}{\bu_{n+1}}+\mb{M}_{n+1}\ddfrac{\ba_{n+1}}{\bu_{n+1}}.
\end{gather}
It thus can be concluded from \eqsref{eq:revised_k} that to account for this displacement dependent damping term, only the stiffness matrix needs to be scaled by the factor $1+\sum\dfrac{m_j\Delta{}t}{2+s_j\Delta{}t}$.